尚须证明复合
\begin{equation}\tag{4.2.9}
\begin{aligned}
&H^\bullet(X,\mathbb Z_\ell)\oplus
 H^{\bullet-2}(\Delta,\mathbb Z_\ell(-1))
 \xrightarrow{\ f^*+i_*g^*\ }
 H^\bullet(\widetilde X,\mathbb Z_\ell)\\[-2pt]
&\hspace{55mm}\xrightarrow{\ f_*\oplus(-g_*i^*)\ }
 H^\bullet(X,\mathbb Z_\ell)\oplus
 H^{\bullet-2}(\Delta,\mathbb Z_\ell(-1))
\end{aligned}
\end{equation}
为恒等映射. 而
\[
f_*f^*=\mathrm{id}\quad\text{在 }H^\bullet(X,\mathbb Z_\ell)
\quad\text{由 (2.2.18)},
\]
\[
g_*i^*f^*=0\quad\text{由 (4.2.5)},
\qquad
f_*i_*g^*=0\quad\text{由 (4.2.7)}.
\]
因此只须证明
\begin{equation}\tag{4.2.10}
-g_*i^*i_*g^*(x)=x
\quad\text{对 }x\in H^\bullet(\Delta,\mathbb Z_\ell(-1)).
\end{equation}

为验证 (4.2.10), 首先注意 $i:\widetilde\Delta\hookrightarrow\widetilde X$ 是一个余维为一的光滑子簇的包含映射. 因而 (SGA 5 VII)
\begin{equation}\tag{4.2.11}
i^*i_*(x)=c_1(\underline N_{\widetilde\Delta/\widetilde X})\,x,
\end{equation}
以及
\begin{equation}\tag{4.2.12}
i_*(x)i_*(y)
=i_*\!\left(xy\,c_1(\underline N_{\widetilde\Delta/\widetilde X})\right).
\end{equation}
记
$\xi_X=\underline{\mathcal O}_X(1)$,
$\xi_\Delta=\underline{\mathcal O}_\Delta(1)$, 以及
$\xi_{\mathbb P^1}=\underline{\mathcal O}_{\mathbb P^1}(1)$.
由 EGA II 8.1.7, 若 $\widetilde{\mathfrak J}$ 为 $\widetilde\Delta$ 在 $\widetilde X$ 中的理想, 则
\begin{equation}\tag{4.2.13}
\begin{aligned}
\underline{\check N}_{\widetilde\Delta/\widetilde X}
&\simeq\widetilde{\mathfrak J}/\widetilde{\mathfrak J}^{,2}
 \simeq\widetilde{\mathfrak J}
       \otimes_{\underline{\mathcal O}_{\widetilde X}}
       \underline{\mathcal O}_{\widetilde\Delta}\\
&\simeq\underline{\mathcal O}_{\widetilde X}(1)
       \otimes_{\underline{\mathcal O}_{\widetilde X}}
       \underline{\mathcal O}_{\widetilde\Delta}
 =\underline{\mathcal O}_{\widetilde\Delta}(1).
\end{aligned}
\end{equation}
这里将 $\widetilde\Delta$ 看作 $\mathbb P(\mathfrak J/\mathfrak J^2)$, 其中 $\mathfrak J$ 为 $\Delta$ 在 $X$ 中的理想. 由于 $\Delta$ 是 $X$ 与两个给定超平面的交, 有典范同构
\[
\mathfrak J/\mathfrak J^2\simeq
\underline{\mathcal O}_\Delta^{,2}(-1),
\]
它诱导
\[
\widetilde\Delta=\mathbb P(\mathfrak J/\mathfrak J^2)
\simeq\mathbb P(\underline{\mathcal O}_\Delta^{,2}(-1))
\xrightarrow[\sim]{\alpha}
\mathbb P(\underline{\mathcal O}_\Delta^{,2})
\simeq\Delta\times\mathbb P^1.
\]
EGA II 4.1.4 的典范同构 $\alpha$ 将重言线丛变为
\[
(\rho i)^*\underline{\mathcal O}_{\mathbb P^1}(1)
\otimes g^*\underline{\mathcal O}_\Delta(-1).
\]
因而 (4.2.13) 成为
\begin{equation}\tag{4.2.14}
\underline{\check N}_{\widetilde\Delta/\widetilde X}
\simeq(\rho i)^*\underline{\mathcal O}_{\mathbb P^1}(1)
       \otimes g^*\underline{\mathcal O}_\Delta(-1).
\end{equation}

\subsection*{引理 4.3.}
\underline{有}
\[
c_1(\underline N_{\widetilde\Delta/\widetilde X})
=g^*c_1(\xi_\Delta)-(\rho i)^*c_1(\xi_{\mathbb P^1}).
\]
现在回到 (4.2.10) 的证明. 由 (4.2.11) 和引理 4.3, 只须证明
\[
g_*\!\left(g^*(x)\left[(\rho i)^*c_1(\xi_{\mathbb P^1})
                       -g^*c_1(\xi_\Delta)\right]\right)=x.
\]
由 (4.2.4), $g_*g^*=0$, 而
\[
\begin{tikzcd}[column sep=huge]
\widetilde\Delta\arrow[r,"g\times\rho_i","\sim"']
  &\Delta\times\mathbb P^1,
\end{tikzcd}
\]
所以
\begin{equation}\tag{4.3.1}
g_*\!\left(g^*(x)(\rho i)^*c_1(\xi_{\mathbb P^1})\right)=x,
\end{equation}
从而
\begin{equation}\tag{4.3.2}
g_*\!\left(g^*(x)c_1(\underline N_{\widetilde\Delta/\widetilde X})\right)
=-x.
\end{equation}

现在计算下式上的代数结构:
\[
H^\bullet(X,\mathbb Z_\ell)\oplus
H^{\bullet-2}(\Delta,\mathbb Z_\ell(-1))
\]
此结构由 (4.2.2) 从杯积传递而来. 设 $a,b$ 属于 $H^\bullet(X,\mathbb Z_\ell)$, $x,y$ 属于 $H^{\bullet-2}(\Delta,\mathbb Z_\ell(-1))$. 先考虑 (4.2.3.1), 并记
\[
(0\oplus x)\wedge(0\oplus y)=A\oplus B.
\]
由 (4.2.12) 和 (4.3.2),
\[
\begin{aligned}
A
&=f_*\!\left((i_*g^*x)(i_*g^*y)\right)\\
&=f_*i_*\!\left(g^*(xy)c_1(\underline N_{\widetilde\Delta/\widetilde X})\right)\\
&=j_*g_*\!\left(g^*(xy)c_1(\underline N_{\widetilde\Delta/\widetilde X})\right)
=-j_*(xy),
\end{aligned}
\]
以及
\[
\begin{aligned}
B
&=-g_*i^*\!\left((i_*g^*x)(i_*g^*y)\right)\\
&=-g_*i^*i_*\!\left(g^*(xy)c_1(\underline N_{\widetilde\Delta/\widetilde X})\right)\\
&=-g_*\!\left(g^*(xy)c_1(\underline N_{\widetilde\Delta/\widetilde X})^2\right).
\end{aligned}
\]
引理 4.3 给出
\[
\begin{aligned}
c_1(\underline N_{\widetilde\Delta/\widetilde X})^2
&=g^*\!\left(c_1(\xi_\Delta)^2\right)
 -2g^*c_1(\xi_\Delta)(\rho i)^*c_1(\xi_{\mathbb P^1}),
\end{aligned}
\]
因为 $c_1(\xi_{\mathbb P^1})^2=0$. 因此, 由 (4.2.4) 和 (4.3.1),
\[
B=0+2xy\,c_1(\xi_\Delta),
\]
这就证明了 (4.2.3.1). 公式 (4.2.3.2) 表达了 $f^*$ 为代数同态这一事实.

为证明 (4.2.3.3), 由投影公式可知 $\operatorname{Ker}f_*$ 在与 $\operatorname{im}f^*$ 的元素相乘下保持稳定. 因而
\[
(a\oplus0)\wedge(0\oplus y)=0\oplus C,
\]
其中
\begin{equation}\tag{4.3.3}
C=-g_*i^*\!\left(f^*(a)i_*g^*(y)\right).
\end{equation}
利用 (4.3.2),
\[
\begin{aligned}
C
&=-g_*\!\left(i^*f^*(a)g^*(y)
      c_1(\underline N_{\widetilde\Delta/\widetilde X})\right)\\
&=-g_*\!\left(g^*j^*(a)g^*(y)
      c_1(\underline N_{\widetilde\Delta/\widetilde X})\right)\\
&=j^*(a)y.
\end{aligned}
\]
这就证明了 (4.2.3.3). 最后, 由杯积的分次交换性, (4.2.3.4) 从 (4.2.3.3) 得出.
